\section{Notation and the line walk}\label{sec:notation}

This section fixes the notation for the line walk and recalls its crossing equation
from~\cite{paperA}. The planar walk itself is set up at the start of Section~\ref{sec:planar}.

\emph{The line walk.} For $a,b\ge1$ let $S_{a,b}(u)=\sum_wu^{\mathrm{sw}(w)}$, the sum
over the words $w$ with $a$ right steps and $b$ left steps, where $\mathrm{sw}(w)$ is the
number of switches of $w$. Each such word has probability
$\frac12p^{N-1}(q/p)^{\mathrm{sw}(w)}$ under the line walk with $N=a+b$ steps, and we say
that it ends in bin $a$. So
$S_{a,b}(q/p)$ is the probability of bin $a$ divided by the probability
$\frac12p^{N-1}$ of an endpoint. For $N\ge2$ let
$S_N=S_{\lfloor N/2\rfloor,\lceil N/2\rceil}$, which belongs to a middle bin.
It is a nonzero polynomial with nonnegative coefficients
and no constant term, so $S_N(u)=1$ has exactly one positive root $u_N$ (the crossing value of
the odds $u=q/p$ of a switch), and
$p_N=1/(1+u_N)$ is the crossing of the middle and the endpoints. Put $L=\log N$,
$z_N=N(1-p_N)$ and $c_{\mathrm A}=\frac12\log(\pi/8)$. By~\cite{paperA}, $z_N$ satisfies the crossing
equation $z_N+\frac12\log z_N=L+c_{\mathrm A}+1/(8z_N)+R_N$ with an explicit remainder
$R_N=O(1/L^2)$.

\emph{Conventions.} The rates $r$ and $s$ are fixed. Unless we say otherwise, every constant,
every implied constant in $O(\cdot)$ and every threshold for $N$ depends only on them. Limits
and $o(\cdot)$ refer to $N\to\infty$ through even integers. For a path of the planar walk we
write $n_j$ for its number of steps in the direction $j$, and we call the vector
$n=(n_0,n_1,n_2,n_3)$ its composition.
